% संपूर्ण मराठी ग्रंथातील प्रकरण 72: सिद्धता-शोध
% हा संपादनयोग्य स्रोतखंड आहे. संकलनासाठी संपूर्ण स्रोतसंग्रहातील
% अक्षररूपे, पूर्वभाग, चिन्हव्याख्या आणि आकृती-स्रोत आवश्यक आहेत.
% मूळ: Open Logic Project; मजकूर आणि रूपांतर CC BY 4.0.
% यंत्रानुवाद, दुरुस्ती आणि तपासणी: OpenAI Codex —
% GPT-5.6 Sol आणि GPT-6 Sol, दोन्ही Ultra effort.
% दुरुस्ती-पुनरावलोकन: GPT-6.1 Sol, Ultra effort.
\chapter{सिद्धता-शोध}\label{pt:ps:chap}










\section{प्रस्तावना}\label{pt:ps:int:sec}

\Log{G3c} किंवा~\Log{N1i} सारख्या सिद्धता
पद्धतींची स्वयंसिद्धके आणि अनुमाननियम ठरवतात की
क्रमवर्ती किंवा सूत्रे यांच्या कोणत्या वृक्षांना
योग्य सिद्धता म्हणायचे. असा क्रमवर्तींचा किंवा
सूत्रांचा वृक्ष दिला असेल (आवश्यक असल्यास
प्रत्येक टप्प्यावर कोणता नियम लागू झाला, किंवा
गृहीतकांच्या आणि त्यांना निरस्त करणाऱ्या अनुमाननियमांच्या
खुणा यांसारखी अतिरिक्त माहितीही असेल), तर
स्वयंसिद्धके आणि नियम यांच्या व्याख्यांवरून तो
वृक्ष योग्य सिद्धता आहे का हे तपासता येते.
पण दिलेल्या क्रमवर्तीची किंवा दिलेल्या गृहीतकांपासून
एखाद्या सूत्राची सिद्धता \emph{शोधणे}
हा वेगळा आणि बहुतेकदा अधिक कठीण प्रश्न आहे.
हाच \emph{सिद्धता-शोध} प्रश्न.

सिद्धता-शोधासाठी एक अगदी सोपी पद्धत आहे.
सूत्रे यांना सांत वर्णमालेतील चिन्हांचे अनुक्रम
समजू शकतो. क्रमवर्ती आणि क्रमवर्तींचे किंवा
सूत्रांचे वृक्षही सूत्रांचे अनुक्रम
मानता येतात (वृक्षाची शाखा दाखवण्यासाठी विशेष
कंस वापरून). चिन्हांचा एखादा अनुक्रम योग्य सिद्धता
दर्शवतो का, हे स्वयंसिद्धके आणि नियम यांमुळे
यांत्रिकरीत्या तपासता येते. परिणामी, सर्व \emph{शक्य}
योग्य सिद्धता यांत्रिकरीत्या निर्माण करता येतात:
सिद्धता दर्शवणाऱ्या वर्णमालेतील चिन्हांचे सर्व
अनुक्रम तयार करून प्रत्येक उमेदवार अनुक्रम योग्य
सिद्धता आहे का ते तपासा. ही सोपी सिद्धता-शोध
पद्धत अशी: दिलेल्या क्रमवर्तीची सिद्धता मिळेपर्यंत
सर्व योग्य सिद्धता निर्माण करा. (किंवा दिलेल्या
गृहीतकांपैकी काही मुक्त गृहीतके असलेल्या दिलेल्या
सूत्राची सिद्धता मिळेपर्यंत तसे करा.)

याहून चांगली पद्धत म्हणजे हाताने सिद्धता
शोधताना वापरायची सहज पद्धत: अंतिम क्रमवर्तीपासून
सुरुवात करा, एक सूत्र निवडा, आणि अनुमाननियम
``उलट्या'' दिशेने लावून एक किंवा अनेक आधारविधाने
निर्माण करा. त्यांची सिद्धता असतील, तर शेवटच्या
अनुमानासह त्यांतून दिलेल्या क्रमवर्तीची सिद्धता
बनते. नैसर्गिक निगमन पद्धतीत सिद्धता
शोधण्यासाठीही अशीच, पण अधिक गुंतागुंतीची, पद्धत
वापरली जाईल.

मात्र ही पद्धत हमखास यशस्वी नाही:
चुकीचा नियम किंवा सूत्रांचा चुकीचा क्रम
निवडल्यास मार्ग बंद होऊ शकतो. ती पूर्णपणे
निर्दिष्टही नाही: प्रत्येक टप्प्यावर निवडण्यासाठी
एकापेक्षा अधिक सूत्र किंवा उलट्या दिशेने
लावण्यासाठी एकापेक्षा अधिक नियम असू शकतात.
\emph{सिद्धता-शोध अल्गोरिदम} म्हणजे पूर्णपणे
निर्दिष्ट अशी प्रक्रिया की सिद्धता अस्तित्वात
असेल तर ती शोधून काढते (म्हणजे ती हमखास यशस्वी आहे).

काही सिद्धता पद्धतींसाठी सिद्धता-शोध
अल्गोरिदम इतरांपेक्षा सोपे असतात. क्रमवर्ती-पद्धतीत
सूत्राचा मुख्य संकारक आणि ते डावीकडे
की उजवीकडे आहे यावर उलट्या दिशेने कोणता नियम
लावायचा हे ठरते. उदाहरणार्थ,~$A \land B$
उजवीकडे असलेल्या~$\Gamma \Sequent \Delta, A \land B$
ची सिद्धता शोधायची असेल, तर~\RightR{\land}
उलटा लावून अनुरूप दोन आधारविधाने
$\Gamma \Sequent \Delta, A$ आणि
$\Gamma \Sequent \Delta, B$ निर्माण करावीत.
पण पद्धत~\Log{G1c} असेल, तर आकुंचनामुळे अधिक
पर्याय मिळतात आणि काम कठीण होते: उलट्या दिशेने
\RightR{\Contraction} लावून
$\Gamma \Sequent \Delta, A \land B, A \land B$
हे आधारविधानही तयार करता येईल. मग
$\Gamma \Sequent \Delta, A \land B, A \land B, A \land B$
हे, आणि असे पुढे. \Log{G2c} मध्ये अडचण अधिक आहे.
\RightR{\land} उलटा लावताना~$\Gamma_1$, $\Gamma_2$
हे~$\Gamma = \Gamma_1 \cup \Gamma_2$ पूर्ण करतील
असा अंदाज, आणि~$\Delta_1$, $\Delta_2$ हे
$\Delta = \Delta_1 \cup \Delta_2$ पूर्ण करतील
असा अंदाज बांधून,~$\Gamma_1 \Sequent \Delta_1, A$
आणि~$\Gamma_2 \Sequent \Delta_2, B$ ही
आधारविधाने वापरून पाहावी लागतात. चुकीच्या अंदाजाने
पुढे मार्ग बंद झाला, तर सुरुवातीला परत जाऊन
वेगळे~$\Gamma_1$, $\Gamma_2$, $\Delta_1$, $\Delta_2$
निवडावे लागतात. सूत्र हे~$A \lor B$
असेल, तर~\RightR{\lor} चे दोन पर्यायांपैकी एक
निवडून~$\Gamma \Sequent \Delta, A$ किंवा
$\Gamma \Sequent \Delta, B$ या दोनपैकी
एक आधारविधान निर्माण करावे लागते. चुकीच्या
निवडीवर पुन्हा मागे फिरावे लागते.

\Log{G3c} मध्ये या अडचणी नाहीत.
तिथे संरचनात्मक नियम नाहीत; मुख्य संकारक
आणि सूत्राची बाजू यांवर अनुमाननियम ठरतो.
म्हणून~\Log{G3c} हे~\Log{G1c} किंवा~\Log{G2c}
पेक्षा सिद्धता-शोधासाठी अधिक सोयीचे आहे.
यशस्वी शोधातून~\Log{G3c} मधील सिद्धता मिळते;
नंतर त्या सिद्धता चे~\Log{G1c} किंवा~\Log{G2c}
(किंवा~\Log{N1c}) मध्ये रूपांतर करता येते.
त्यामुळे~\Log{G3c} साठी सिद्धता-शोध अल्गोरिदम
आणि~\Log{G3c} सिद्धता चे इतर पद्धतींतील
सिद्धता मध्ये रूपांतर करणारे अल्गोरिदम मिळाले,
तर त्या इतर पद्धतींचा सिद्धता-शोध प्रश्नही सुटतो.

सिद्धता-शोध अल्गोरिदम हमखास यशस्वी आहे
हे कसे समजेल? त्याने सिद्धता शोधली नाही,
तर सिद्धता अस्तित्वातच नाही हे दाखवावे लागेल.
कारण अल्गोरिदम चुकीचा निवडला, तर सिद्धता
असूनही ती न सापडण्याचे प्रसंग येऊ शकतात.
अगदी सोप्या पद्धतीत ही अडचण नसते: ती सर्व
शक्य सिद्धता तपासते. पण सिद्धता-शोध
अल्गोरिदम कार्यक्षम असावा आणि शक्य तितके कमी
प्रयत्न वापरावेत अशी अपेक्षा असते.

पुढील बंद-पद शोध आणि त्याच्या पद-प्रतिमानाची सिद्धता
वाक्यांच्या क्रमवर्तींसाठी मांडली आहे; खालील पूर्तिचे
वर्णन या व्याप्तीत वाचायचे.
सिद्धता-शोध अल्गोरिदम हमखास यशस्वी आहे
हे दाखवण्याची मुख्य कल्पना अशी: शोध अयशस्वी
झाल्यास त्याच्या अपयशाच्या स्वरूपावरून त्या
क्रमवर्तीची सिद्धता अस्तित्वात असू शकत नाही हे दाखवा.
हे दाखवण्यासाठी ती वैध नाही असे दाखवू.
अभिजात प्रथम-क्रम तर्कशास्त्रात
$\Gamma \Sequent \Delta$ ही क्रमवर्ती वैध
असण्याचा अर्थ असा की प्रत्येक संरचना
मध्ये~$\Gamma$ मधील किमान एक सूत्र
असत्य असेल किंवा~$\Delta$ मधील किमान एक
सूत्र सत्य असेल. \Log{G3c} \emph{निर्दोष}
आहे; म्हणजे तो फक्त वैध क्रमवर्ती सिद्ध करतो.
हे सिद्धता वर विगमन करून सिद्ध होते:
स्वयंसिद्धके स्पष्टपणे वैध आहेत आणि नियमांची
आधारविधाने वैध असतील तर निष्कर्षही वैध असतो.
सिद्धता-शोध अल्गोरिदम हमखास यशस्वी आहे हे
दाखवण्यासाठी,~$\Gamma \Sequent \Delta$ ची
सिद्धता शोधण्यात अपयश आल्यास अशी
संरचना परिभाषित करता येते हे दाखवतो
की तिच्यात~$\Gamma$ मधील सर्व सूत्रे
सत्य आणि~$\Delta$ मधील सर्व सूत्रे
असत्य आहेत. यावरून~\Log{G3c} \emph{संपूर्ण}
आहे हेही सिद्ध होते:~$\Gamma \Sequent \Delta$
वैध असेल, तर तिची~\Log{G3c}-सिद्धता असते.
अयशस्वी सिद्धता-शोधामुळे~$\Gamma \Sequent \Delta$
अवैध ठरते; म्हणून वैध क्रमवर्तीसाठीचा प्रत्येक
सिद्धता-शोध यशस्वी ठरलाच पाहिजे.













\section{\Log{G3c} साठी सिद्धता-शोध अल्गोरिदम}\label{pt:ps:sal:sec}

\Log{G3c} साठी सिद्धता-शोध अल्गोरिदम मांडू.
हा एकमेव शक्य अल्गोरिदम नाही आणि सर्वांत कार्यक्षमही
नाही. पण तो बरोबर आहे:~$\Gamma \Sequent \Delta$
या क्रमवर्तीची सिद्धता असेल, तर अखेरीस तो
थांबतो आणि ती सिद्धता शोधतो. तो विश्वासार्हही
आहे: सिद्धता न सापडता तो थांबला, किंवा कधीच
समाप्त झाला नाही, तर~$\Gamma \Sequent \Delta$
मुळातच वैध नाही.
या बंद-पद शोधात सुरुवातीच्या क्रमवर्तीतील सर्व सूत्रे
वाक्ये, म्हणजे मुक्त चल नसलेली सूत्रे, असावीत.

अल्गोरिदमचा अत्यावश्यक गुणधर्म असा की
$\Gamma \Sequent \Delta$ मधील किंवा सिद्धता-शोधात
निर्माण झालेल्या क्रमवर्तीतील प्रत्येक सूत्र
न्यूनीकृत होते: उलट्या दिशेने लावलेल्या
अनुमाननियमात त्याला मुख्य सूत्र मानले जाते,
आणि गरजेइतक्या वेळा असे केले जाते.
या गुणधर्माला \emph{न्याय्यता} म्हणू.

अल्गोरिदम टप्प्याटप्प्याने चालतो.
प्रत्येक टप्प्याचे फलित क्रमवर्तींचा वृक्ष असतो.
पुढील टप्प्याचे फलित मिळवताना आधीच स्वयंसिद्धक
नसलेल्या प्रत्येक सर्वांत वरच्या क्रमवर्तीतील
सूत्राचे न्यूनीकरण करतो.

~$0$ व्या टप्प्यावर~$\Gamma \Sequent \Delta$
ने सुरुवात करतो. न्याय्यता सुनिश्चित करण्यासाठी
$\Gamma$ आणि~$\Delta$ मधील सूत्रांच्या
प्रत्येक आस्थितीला, म्हणजे वेगवेगळ्या सूत्रे
च्या आस्थितींना, वेगळा संख्यात्मक निर्देशांक
देतो. उदाहरणार्थ, सूत्रे डावीकडून उजवीकडे
मोजा. या संख्या ऊर्ध्वनिर्देशांक म्हणून लिहितो.
उदाहरणार्थ,~$A, B \Sequent C$ ची सिद्धता
शोधायची असेल, तर~$0$ व्या टप्प्याचे फलित
$A^1, B^2 \Sequent C^3$ असते. अशा निर्देशांकित
क्रमवर्तीतील ऊर्ध्वनिर्देशांकांपैकी सर्वांत मोठा
~$n$ हा तिचा \emph{कमाल निर्देशांक}
(उदाहरणात~$3$) आहे.

क्रमवर्तीत संख्यापक असतील, तर
उजवीकडील~$\lexists[x][A(x)]$ किंवा डावीकडील
$\lforall[x][A(x)]$ न्यूनीकृत करताना कोणती पदे
वापरायची तेही ठरवावे लागते. ~$\Gamma$ आणि~$\Delta$
मधील फलनचिन्हे वापरून~स्थिरांक
$\Obj c_0$, $\Obj c_1$, $\Obj c_2$, \dots
यांपासून बनणारी पदेच लागतात. या स्थिरांक
पासून बनणाऱ्या सर्व पदांचा
एक निश्चित क्रम दिला आहे असे मानू; म्हणजे,
उदाहरणार्थ,~``$\Gamma \Sequent \Delta$ मध्ये
न आलेला पहिला स्थिरांक'' असे म्हणता येईल.
अल्गोरिदमने बनवलेल्या वृक्षातील प्रत्येक
निर्देशांकित क्रमवर्तीशी स्थिरांकाचा एक संच
संबंधित असतो. ~$0$ व्या टप्प्यातील क्रमवर्तीशी
तिच्यात येणाऱ्या सर्व स्थिरांकाचा संच
(असतील तर) संबंधित असतो; त्यात एकही
स्थिरांक नसतील, तर~$\{\Obj c_0\}$ असतो.

~$n$ व्या टप्प्यावर संबंधित स्थिरांक
च्या संचांसह निर्देशांकित क्रमवर्तींचा वृक्ष
अल्गोरिदमने निर्माण केला आहे असे मानू.
~$n+1$ व्या टप्प्यावर प्रत्येक सर्वांत वरच्या
क्रमवर्तीसाठी पुढीलपैकी योग्य ती कृती करतो.

~$A$ हे सूत्र क्रमवर्तीच्या डाव्या
आणि उजव्या दोन्ही बाजूंना असेल, किंवा डावीकडे
$\lfalse$ असेल, तर काही करत नाही: अशा क्रमवर्तींच्या
~\Log{G3c} मधील सिद्धता असतात; या विशिष्ट
सर्वांत वरच्या क्रमवर्तीसाठी पुढे काही करायचे नसते.

क्रमवर्तीत अणुरूप नसलेले एकही सूत्र
नसेल, तर अल्गोरिदम अयशस्वी म्हणून थांबवा.
$\Gamma \Sequent \Delta$ ची सिद्धता नाही.

अन्यथा सर्वांत लहान निर्देशांक~$i$ असलेले
अणुरूप नसलेले सूत्र~$A$ निवडा. मग विचाराधीन
सर्वांत वरची क्रमवर्ती~$A^i, \Pi \Sequent \Lambda$
किंवा~$\Pi \Sequent \Lambda, A^i$ अशी असते.
त्या क्रमवर्तीचा कमाल निर्देशांक~$k$ आणि
तिच्याशी संबंधित स्थिरांकाचा संच~$C$ असू द्या.

~$A^i$ चे ``न्यूनीकरण'' करू:~$A$
चा मुख्य संकारक आणि~$A^i$ डावीकडे आहे
की उजवीकडे, यांवर ठरणाऱ्या अनुमाननियमानुसार
नव्या सर्वांत वरच्या क्रमवर्ती तयार करतो.

\begin{enumerate}
  \item $A \ident B \land C$ असून ते डावीकडे
  असेल, तर~$(B \land C)^i, \Pi \Sequent \Lambda$
  च्या वर एक नवी सर्वोच्च क्रमवर्ती जोडा:
  \[
  \Axiom$B^{k+1}, C^{k+2}, \Pi \fCenter \Lambda$
  \RightLabel{\LeftR{\land}}
  \UnaryInf$(B \land C)^i, \Pi \fCenter \Lambda$
  \DisplayProof
  \]
  नव्या क्रमवर्तीशी संबंधित स्थिरांकाचा संच~$C$ आहे.
  \item $A \ident B \land C$ असून ते उजवीकडे
  असेल, तर~$\Pi \Sequent \Lambda, (B \land C)^i$
  च्या वर दोन नव्या सर्वोच्च क्रमवर्ती जोडा:
  \[
  \Axiom$\Pi \fCenter \Lambda, B^{k+1}$
  \Axiom$\Pi \fCenter \Lambda, C^{k+1}$
  \RightLabel{\RightR{\land}}
  \BinaryInf$\Pi \fCenter \Lambda, (B \land C)^i$
  \DisplayProof
  \]
  प्रत्येक नव्या क्रमवर्तीशी संबंधित स्थिरांकाचा संच~$C$ आहे.

  \item $A \ident \lnot B$,
  $A \ident B \lor C$ आणि~$A \ident B \lif C$
  हे प्रसंगही याच प्रकारे हाताळता येतात;
  ते सरावासाठी ठेवले आहेत.

  \item $A \ident \lexists[x][B(x)]$ असून ते
  डावीकडे असेल, तर~$\lexists[x][B(x)]^i,
  \Pi \Sequent \Lambda$ च्या वर एक नवी
  सर्वोच्च क्रमवर्ती जोडा:
  \[
  \Axiom$B(c)^{k+1}, \Pi \fCenter \Lambda$
  \RightLabel{\LeftR{\lexists}}
  \UnaryInf$\lexists[x][B(x)]^i, \Pi \fCenter \Lambda$
  \DisplayProof
  \]
  येथे~$c$ हा~$C$ मध्ये नसलेला पहिला स्थिरांक
  आहे. नव्या क्रमवर्तीशी संबंधित स्थिरांक
  चा संच~$C \cup \{c\}$ आहे.
  \item $A \ident \lexists[x][B(x)]$ असून ते
  उजवीकडे असेल, तर~$\Pi \Sequent \Lambda,
  \lexists[x][B(x)]^i$ च्या वर एक नवी
  सर्वोच्च क्रमवर्ती जोडा:
  \[
  \Axiom$\Pi \fCenter \Lambda, \lexists[x][B(x)]^{k+1}, B(t)^{k+2}$
  \RightLabel{\RightR{\lexists}}
  \UnaryInf$\Pi \fCenter \Lambda, \lexists[x][B(x)]^i$
  \DisplayProof
  \]
  येथे~$t$ हे~$C$ मधीलच स्थिरांक असलेले,
  आणि या क्रमवर्तीखाली उजवीकडील
  $\lexists[x][B(x)]$ च्या आधीच्या न्यूनीकरणात
  न वापरलेले पहिले पद आहे. अशी सर्व पदे वापरून
  झाली असतील, तर~$C$ मधील पहिला स्थिरांक पुन्हा वापरा.
  $C$ सुरुवातीपासून अरिक्त आहे आणि पुढे त्यात फक्त भर
  पडते; म्हणून हा स्थिरांक उपलब्ध आहे आणि~$C$ बाहेरचा
  नवा स्थिरांक नकळत क्रमवर्तीत येत नाही.
  नव्या क्रमवर्तीशी संबंधित स्थिरांकाचा संच~$C$.
  \item $A \ident \lforall[x][B(x)]$
  हे प्रसंगही याच प्रकारे हाताळता येतात;
  ते सरावासाठी ठेवले आहेत.
\end{enumerate}

~$n+1$ व्या टप्प्यावर कोणत्याही सर्वांत वरच्या
क्रमवर्तीच्या वर नवी क्रमवर्ती जोडली गेली नसेल,
तर अल्गोरिदम यशस्वीपणे संपतो. अशा वेळी
~$n+1$ व्या टप्प्यातील वृक्षाचे सर्व निर्देशांक
काढल्यावर~$\Gamma \Sequent \Delta$ ची
सिद्धता मिळते. सर्वांत वरच्या क्रमवर्तींचे
रूप~$A, \Pi \Sequent \Lambda, A$ किंवा
$\lfalse, \Pi \Sequent \Lambda$ असे असते.
सर्व अनुमाने~\Log{G3c} मधील योग्य अनुमाने असतात.
विधानीय अनुमानांसाठी हे स्पष्ट आहे. प्रत्येक
टप्प्यावर क्रमवर्तीशी संबंधित स्थिरांकाच्या
संच~$C$ मध्ये त्या क्रमवर्तीतील सर्व स्थिरांक असतात.
म्हणून~\LeftR{\lexists} आणि~\RightR{\lforall}
यांनी न्यूनीकरण करताना आयगेन-चराची अट पूर्ण होते:
आयगेन-चर~$c$ हा निष्कर्षाशी संबंधित संच~$C$
मध्ये नसलेला स्थिरांक असल्याने तो
~$c$ निष्कर्षातही आलेला नव्हता. म्हणून:

\begin{prop}
~\Log{G3c} साठीचा सिद्धता-शोध अल्गोरिदम बरोबर आहे:
तो यशस्वीपणे संपला, तर सुरुवातीच्या
क्रमवर्ती~$\Gamma \Sequent \Delta$ ची
सिद्धता असते.
\end{prop}













\section{संपूर्णता}\label{pt:ps:cpl:sec}

आता सिद्धता-शोध अल्गोरिदम विश्वासार्ह आहे हे दाखवू:
तो सिद्धता न सापडता अयशस्वी म्हणून थांबला तरी किंवा कधीच समाप्त झाला नाही तरी,
सुरुवातीच्या क्रमवर्ती~$\Gamma \Sequent \Delta$ ची
सिद्धता नसते. यासाठी अशी संरचना~$\Struct M$
बांधू की प्रत्येक~$A \in \Gamma$ साठी~$\Sat{M}{A}$
आणि प्रत्येक~$A \in \Delta$ साठी~$\Sat/{M}{A}$.
म्हणजे~$\Gamma$ मधील सर्व सूत्रे सत्य,
तर~$\Delta$ मधील सर्व सूत्रे हे~$\Struct M$
मध्ये असत्य आहेत. त्यामुळे~$\Gamma \Sequent \Delta$
वैध नाही. \Log{G3c} निर्दोष असल्याने
$\Gamma \Sequent \Delta$ ची सिद्धता नाही.
येथे सुरुवातीच्या क्रमवर्तीतील सर्व सूत्रे वाक्ये आहेत,
म्हणजे त्यांत मुक्त चल नाहीत. पुढील बंद-पद अल्गोरिदम
आणि पद-प्रतिमानाची सिद्धता या व्याप्तीत वापरायची.

सूत्रांची जोडी~$\Theta, \Xi$ आणि
स्थिरांकाचा संच~$C$ यांपासून~$\Struct{M}$ बांधतो.
हे~$\Theta, \Xi$ आणि~$C$ आपल्याला अयशस्वी
(अयशस्वी म्हणून थांबलेल्या किंवा समाप्त न होणाऱ्या) सिद्धता-शोधाच्या
\emph{अपयश-शाखे}तून मिळतात. अपयश-शाखा म्हणजे
क्रमवर्तींचा~$\Pi_n \Sequent \Lambda_n$ अनुक्रम
आणि~स्थिरांकाचे संच~$C_n$; येथे
$\Pi_0 \Sequent \Lambda_0$ ही सुरुवातीची
अंतिम क्रमवर्ती~$\Gamma \Sequent \Delta$ असते,
आणि~$\Pi_{n+1} \Sequent \Lambda_{n+1}$
हे~$n+1$ व्या टप्प्यावर~$\Pi_n \Sequent \Lambda_n$
पासून निर्माण झालेले असे आधारविधान असते की
अल्गोरिदमला त्याची सिद्धता सापडत नाही.
प्रत्येक~$n$ साठी असे आधारविधान असलेच पाहिजे;
नसते तर अल्गोरिदमला सिद्धता सापडली असती.
~$C_n$ हा~$\Pi_n \Sequent \Lambda_n$
शी संबंधित स्थिरांकाचा संच घेऊ.

अल्गोरिदम फक्त अणुरूप सूत्रे असलेल्या
सर्वांत वरच्या क्रमवर्तीवर अयशस्वी म्हणून थांबला, तर त्या
क्रमवर्तीवर संपणारी शाखा ही अपयश-शाखा आहे.
अल्गोरिदम कधीच समाप्त झाला नाही, तर तो वाढत्या
मोठ्या वृक्षांची निर्मिती करत राहतो. हे वृक्ष
वाढून एक अनंत वृक्ष तयार होतो असे समजू शकतो
(अल्गोरिदम अनंत टप्पे चालल्यानंतर जो वृक्ष असेल तो).
हा अनंत वृक्ष सांत-शाखित असतो: प्रत्येक शिखराला
फक्त एक किंवा दोन बालशिखरे असतात---म्हणजे त्याच्या
तत्काळ वरच्या क्रमवर्ती. क्योनिग (K\H{o}nig)
यांच्या उपप्रमेयानुसार,
प्रत्येक अनंत सांत-शाखित वृक्षाला अनंत शाखा असते.
सिद्धता-शोध कायम चालू राहिल्यास अशी अनंत शाखा
अपयश-शाखा ठरते.

~$\Pi_n \Sequent \Lambda_n$ आणि~$C_n$
यांचा अनुक्रम अपयश-शाखा असेल, तर
$\Theta = \bigcup_n \Pi_n$ आणि~$\Xi = \bigcup_n \Lambda_n$
घेऊ (सूत्रे वरील निर्देशांक आणि त्यांच्या
अनेक आस्थिती काढून); तसेच~$C = \bigcup_n C_n$ घेऊ.

आता~$\Struct{M}$ परिभाषित करू.
\begin{enumerate}
\item $\Domain{M}$ हा~$C$ मधील स्थिरांक आणि
$\Gamma \Sequent \Delta$ मधील फलने
(असतील तर) वापरून, पण चले न वापरता,
बनवता येणाऱ्या पदांचा संच आहे.
\item $c \in C$ असेल, तर~$\Assign{c}{M} = c$.
\item $f$ हे~$\Gamma \Sequent \Delta$ मधील
$m$-स्थानी फलन असेल आणि~$t_1$,
\dots, $t_m \in \Domain{M}$ असतील, तर
$\Assign{f}{M}(t_1, \dots, t_m) = f(t_1, \dots, t_m)$.
\item $R$ हे~$\Gamma \Sequent \Delta$ मधील
$m$-स्थानी विधेय असेल, तर
$\Assign{R}{M} = \Setabs{\tuple{t_1, \dots, t_m} \in
\Domain{M}^m}{R(t_1, \dots, t_m) \in \Theta}$.
\end{enumerate}
~$\Struct{M}$ हे \emph{पद-प्रतिमान} आहे, कारण
त्याचे विश्व पदांचा संच आहे आणि स्थिरांक व
फलने यांचा अर्थ असा लावला आहे की
$\Value{t}{M} = t$ मिळते. ~$\Theta$ मधील अणुरूप
सूत्रे वापरून~$\Assign{R}{M}$ अशी परिभाषित
केली आहे की~$\Sat{M}{R(t_1, \dots, t_m)}$ तेव्हाच
आणि तेव्हाच जेव्हा~$R(t_1, \dots, t_m) \in \Theta$.

\begin{lem}\label{pt:ps:cpl:lem:truth}
प्रत्येक~$A \in \Theta \cup \Xi$ साठी:
\begin{enumerate}
  \item $A \in \Theta$ असेल, तर~$\Sat{M}{A}$.
  \item $A \in \Xi$ असेल, तर~$\Sat/{M}{A}$.
\end{enumerate}
\end{lem}

\begin{proof}
~$\depth{A}$ वर विगमन करू.

प्रथम~$A$ अणुरूप आहे असे मानू;
म्हणजे~$A \ident \lfalse$ किंवा
$A \ident R(t_1, \dots, t_m)$.

~$\Pi_n \Sequent \Lambda_n$ ही अपयश-शाखा
असल्याने कोणत्याही~$\Pi_n$ मध्ये~$\lfalse$
असू शकत नाही (नाहीतर~$\Pi_n \Sequent \Lambda_n$
स्वयंसिद्धक ठरेल). ~$\Struct{M}$ च्या व्याख्येनुसार
$\Sat{M}{R(t_1, \dots, t_m)}$ तेव्हा आणि केवळ तेव्हा
जेव्हा~$R(t_1, \dots, t_m) \in \Theta$.
दोन्ही मिळून: $A \in \Theta$ असेल, तर
$\Sat{M}{A}$.

~$A$ अणुरूप असून~$A \in \Pi_n$ असेल,
तर~$A \in \Pi_{n+1}$; आणि~$A \in \Lambda_n$
असेल, तर~$A \in \Lambda_{n+1}$.
(सिद्धता-शोध अल्गोरिदम पुढील क्रमवर्ती कशा
बनवतो यावरून हे मिळते.) म्हणून~$A \in \Theta \cap \Xi$
असेल, तर एखाद्या~$n$ साठी~$A \in \Pi_n \cap
\Lambda_n$. मग~$\Pi_n \Sequent \Lambda_n$
स्वयंसिद्धक ठरते; पण अपयश-शाखेत स्वयंसिद्धके
असू शकत नाहीत. म्हणून अणुरूप~$A$ साठी
$A \notin \Theta \cap \Xi$; परिणामी~$A \in \Xi$
असेल तर~$A \notin \Theta$. त्यामुळे~$\Struct{M}$
च्या व्याख्येनुसार~$A \in \Xi$ असताना
$\Sat/{M}{A}$.

विगमन-टप्प्यासाठी,~$A$ पेक्षा कमी
खोली असलेल्या सर्व सूत्रे साठी
(1) आणि~(2) खरे आहेत असे मानू. ~$A$ साठी
(1) व~(2) सिद्ध करण्यासाठी प्रसंग वेगळे पाहू.

येथील $A$ अणुरूप नसलेले आहे. या परिच्छेदात
“सर्वांत लहान निर्देशांक” म्हणजे अणुरूप नसलेल्या
आस्थितींमधील सर्वांत लहान निर्देशांक; अणुरूप आस्थिती
जपल्या जातात, पण न्यूनीकरणासाठी निवडल्या जात नाहीत.
प्रथम,~$\Pi_n \Sequent \Lambda_n$ मध्ये
$A^i$ आल्यास ते~$\Pi_{n+1} \Sequent \Lambda_{n+1}$
मध्येही~$A^i$ म्हणून येते, जोपर्यंत~$i$ हा~$\Pi_n \Sequent \Lambda_n$
मधील सर्वांत लहान निर्देशांक नसतो. प्रत्येक
टप्प्यावर जोडलेल्या सर्वांत वरच्या क्रमवर्तींत
सर्वांत लहान निर्देशांक वाढतो. म्हणून अखेरीस
$A^i$ असलेली अशी~$\Pi_n \Sequent \Lambda_n$
मिळते की~$i$ हा सर्वांत लहान निर्देशांक असतो.
~$n+1$ व्या टप्प्यावर अल्गोरिदम~$A^i$ चे
न्यूनीकरण करतो. दुसऱ्या शब्दांत,
$A \in \Theta$ असेल, तर एखाद्या~$n$ साठी
$\Pi_n = \Pi_n', A^i$ आणि~$i$ सर्वांत
लहान निर्देशांक असतो. तसेच~$A \in \Xi$
असेल, तर एखाद्या~$n$ साठी
$\Lambda_n = \Lambda_n', A^i$ आणि~$i$
सर्वांत लहान निर्देशांक असतो.

\begin{enumerate}
  \item $A \in \Theta$ आणि~$A \ident B \land C$
  असेल, तर एखाद्या~$n$ साठी
  $\Pi_n = \Pi_n', (B \land C)^i$.
  ~$\Pi_{n+1} \Sequent \Lambda_{n+1}$ हे
  \LeftR{\land} चे अनुरूप आधारविधान आहे;
  म्हणजे~$\Pi_{n+1} = \Pi_n', B^{k+1}, C^{k+2}$.
  त्यामुळे~$B, C \in \Theta$. विगमन गृहीतकानुसार
  $\Sat{M}{B}$ आणि~$\Sat{M}{C}$.
  ~$\Sat{M}{}$ च्या व्याख्येवरून
  $\Sat{M}{B \land C}$.

  \item $A \in \Xi$ आणि~$A \ident B \land C$
  असेल, तर एखाद्या~$n$ साठी
  $\Lambda_n = \Lambda_n', B \land C$.
  ~$\Pi_{n+1} \Sequent \Lambda_{n+1}$ हे
  निष्कर्ष~$\Pi_n \Sequent \Lambda'_n, B \land C$
  असलेल्या~\RightR{\land} च्या दोन आधारविधानांपैकी
  एक आहे; म्हणजे~$\Lambda_{n+1} = \Lambda_n', B^{k+1}$
  किंवा~$\Lambda_{n+1} = \Lambda_n', C^{k+1}$.
  त्यामुळे~$B \in \Xi$ किंवा~$C \in \Xi$.
  विगमन गृहीतकानुसार~$\Sat/{M}{B}$ किंवा
  $\Sat/{M}{C}$. ~$\Sat{M}{}$ च्या व्याख्येवरून
  $\Sat/{M}{B \land C}$.

  \item $A \ident \lnot B$,
  $A \ident B \lor C$ आणि~$A \ident B \lif C$
  हे प्रसंगही याच प्रकारे हाताळता येतात;
  ते सरावासाठी ठेवले आहेत.

  \item $A \in \Theta$ आणि
  $A \ident \lexists[x][B(x)]$ असेल,
  तर एखाद्या~$n$ साठी
  $\Pi_n = \Pi_n', \lexists[x][B(x)]^i$.
  ~$\Pi_{n+1} \Sequent \Lambda_{n+1}$ हे
  \LeftR{\lexists} चे अनुरूप आधारविधान आहे;
  म्हणजे एखाद्या~$c$ साठी
  $\Pi_{n+1} = \Pi_n', B(c)^{k+1}$.
  त्यामुळे~$B(c) \in \Theta$ आणि~$c \in C$.
  विगमन गृहीतकानुसार~$\Sat{M}{B(c)}$.
  ~$\Sat{M}{}$ च्या व्याख्येवरून
  $\Sat{M}{\lexists[x][B(x)]}$.

  \item $A \in \Xi$ आणि
  $A \ident \lexists[x][B(x)]$ असेल,
  तर एखाद्या~$n$ साठी
  $\Lambda_n = \Lambda_n', \lexists[x][B(x)]^i$.
  प्रत्येक वेळी~$\lexists[x][B(x)]$ चे
  न्यूनीकरण करताना~$\lexists[x][B(x)]$
  आधारविधानाच्या उजव्या
  बाजूस नव्या निर्देशांकासह राहते. म्हणून
  अपयश-शाखेत ते उजव्या बाजूस असेल, तर
  $\lexists[x][B(x)]$ चे अनंत वेळा
  न्यूनीकरण होते. प्रत्येक~$t \in \Domain{M}$
  हे अखेरीस अपयश-शाखेत~``$C_n$ मधीलच स्थिरांक
  असलेले आणि उजवीकडील~$\lexists[x][B(x)]$
  च्या आधीच्या कोणत्याही न्यूनीकरणात न वापरलेले
  पहिले पद'' ठरते. त्यामुळे प्रत्येक
  $t \in \Domain{M}$ साठी एखाद्या~$n$ ला
  $B(t) \in \Lambda_n$, आणि म्हणून~$B(t) \in \Xi$.
  विगमन गृहीतकानुसार प्रत्येक~$t \in \Domain{M}$
  साठी~$\Sat/{M}{B(t)}$. ~$\Sat{M}{}$
  च्या व्याख्येवरून~$\Sat/{M}{\lexists[x][B(x)]}$.

  \item $A \ident \lforall[x][B(x)]$
  हे प्रसंगही याच प्रकारे हाताळता येतात;
  ते सरावासाठी ठेवले आहेत.
\end{enumerate}
\end{proof}

\begin{cor}\label{pt:ps:cpl:cor:complete}
~\Log{G3c} साठीचा सिद्धता-शोध अल्गोरिदम संपूर्ण
आहे: तो अयशस्वी म्हणून थांबला किंवा कधीच समाप्त झाला नाही,
तर सुरुवातीची क्रमवर्ती~$\Gamma \Sequent \Delta$
वैध नसते आणि म्हणून तिची सिद्धता नसते.
\end{cor}

\begin{proof}
  अल्गोरिदम अयशस्वी म्हणून थांबला किंवा कधीच समाप्त झाला नाही,
  तर अशी अपयश-शाखा असते की तिच्यापासून
  $\Struct{M}$ परिभाषित करता येते.
  \ref{pt:ps:cpl:lem:truth} नुसार प्रत्येक~$A \in \Gamma$
  साठी~$\Sat{M}{A}$ आणि प्रत्येक~$A \in \Delta$
  साठी~$\Sat/{M}{A}$. (लक्षात घ्या की
  $\Pi_0 = \Gamma$ आणि~$\Lambda_0 = \Delta$;
  म्हणून~$\Gamma \subseteq \Theta$ आणि
  $\Delta \subseteq \Xi$.) त्यामुळे
  $\Gamma \Sequent \Delta$ वैध नाही.
  ~\Log{G3c} च्या निर्दोषतेवरून
  $\Gamma \Sequent \Delta$ ची सिद्धता नाही.
\end{proof}

\begin{cor}
वाक्यांच्या क्रमवर्तींसाठी~\Log{G3c} संपूर्ण आहे:
$\Gamma \Sequent \Delta$ वैध असेल,
तर~$\Log{G3c} \Proves \Gamma \Sequent \Delta$.
\end{cor}

\begin{proof}
  व्यत्यस्त विधान सिद्ध करू.
  $\Log{G3c} \Proves/ \Gamma \Sequent \Delta$
  असे मानू. मग शोधण्याजोगी सिद्धता नसल्याने
  सिद्धता-शोध अल्गोरिदम अयशस्वी म्हणून थांबला पाहिजे
  किंवा कायम चालू राहिला पाहिजे.
  \ref{pt:ps:cpl:cor:complete} नुसार
  $\Gamma \Sequent \Delta$ वैध नाही.
\end{proof}













\section{टॅब्लो}\label{pt:ps:tab:sec}

टॅब्लो पद्धती या क्रमवर्ती-कलनाशी जवळून
संबंधित सिद्धता पद्धती आहेत. हाताने किंवा
संगणकीय अंमलबजावणीत सिद्धता-शोधासाठी त्या विशेष
सोयीच्या आहेत. टॅब्लोत क्रमवर्तींचे वृक्ष न बांधता
सिद्धता ही सूत्रांचा वृक्ष असते.
पण नैसर्गिक निगमनापेक्षा वेगळे म्हणजे त्यांचे नियम
क्रमवर्ती-कलनाच्या नियमांशी जुळतात. टॅब्लो
सिद्धता म्हणजे मूलतः प्रतिमानाचा स्पष्ट शोध.
म्हणून त्यांना कधी \emph{चिन्हार्थविषयक} टॅब्लो,
तर कधी \emph{सत्यता-वृक्ष} म्हणतात.

\emph{चिन्हांकित} टॅब्लो हा चिन्हांकित
सूत्रांचा वृक्ष आहे. चिन्हांकित सूत्राचे रूप
$\sFmla{\True}{A}$ किंवा~$\sFmla{\False}{A}$
असे असू शकते. वृक्ष खालच्या दिशेने वाढतात.
सुरुवातीला आपल्याला काय सिद्ध करायचे ते
दर्शवणारी सूत्रे असतात. उदाहरणार्थ,
$A, B \Sequent C, D$ सिद्ध करण्यासाठी
वृक्षाची सुरुवात
$\sFmla{\True}{A}, \sFmla{\True}{B},
\sFmla{\False}{C}, \sFmla{\False}{D}$ अशी
होईल. ~$A$ आणि~$B$ सत्य, पण~$C$ आणि~$D$
असत्य ठरवणारी संरचना म्हणजे
$A, B \Sequent C, D$ अवैध असल्याचे
दाखवणारी संरचना.

टॅब्लोचे तळाचे पानशिखर शाखा-विस्तार
नियमांनी वाढवता येते. हे नियम त्याच शाखेत नवी
शिखरे जोडतात किंवा खाली दोन भावंड शिखरे जोडून
शाखा दोन नव्या शाखांत विभागतात. त्याच शाखेतील
त्या शिखराच्या वरच्या कोणत्याही चिन्हांकित
सूत्राला नियम लागू करता येतात. अभिजात
टॅब्लो पद्धत~\Log{Tc} साठीचे शाखा-विस्तार नियम
\ref{pt:ps:tab:tab:Tc} मध्ये दिले आहेत.








\begin{table}
\begin{center}
\begin{tabular}{@{}cc@{}}
\AxiomC{\sFmla{\True}{\lnot A}}
\RightLabel{\TRule{\True}{\lnot}}
\UnaryInfC{\sFmla{\False}{A}}
\DisplayProof
&
\AxiomC{\sFmla{\False}{\lnot A}}
\RightLabel{\TRule{\False}{\lnot}}
\UnaryInfC{\sFmla{\True}{A}}
\DisplayProof
\\[6ex]
\AxiomC{\sFmla{\True}{A \land B}}
\RightLabel{\TRule{\True}{\land}}
\UnaryInfC{\sFmla{\True}{A}}
\noLine
\UnaryInfC{\sFmla{\True}{B}}
\DisplayProof
&
\AxiomC{\sFmla{\False}{A \land B}}
\RightLabel{\TRule{\False}{\land}}
\UnaryInfC{$\sFmla{\False}{A} \quad \mid \quad \sFmla{\False}{B}$}
\DisplayProof
\\[6ex]
\AxiomC{\sFmla{\True}{A \lor B}}
\RightLabel{\TRule{\True}{\lor}}
\UnaryInfC{$\sFmla{\True}{A} \quad \mid \quad \sFmla{\True}{B}$}
\DisplayProof
&
\AxiomC{\sFmla{\False}{A \lor B}}
\RightLabel{\TRule{\False}{\lor}}
\UnaryInfC{\sFmla{\False}{A}}
\noLine
\UnaryInfC{\sFmla{\False}{B}}
\DisplayProof
\\[6ex]
\AxiomC{\sFmla{\True}{A \lif B}}
\RightLabel{\TRule{\True}{\lif}}
\UnaryInfC{$\sFmla{\False}{A} \quad \mid \quad \sFmla{\True}{B}$}
\DisplayProof
&
\AxiomC{\sFmla{\False}{A \lif B}}
\RightLabel{\TRule{\False}{\lif}}
\UnaryInfC{\sFmla{\True}{A}}
\noLine
\UnaryInfC{\sFmla{\False}{B}}
\DisplayProof
\\[6ex]
\AxiomC{\sFmla{\True}{\lforall[x][A(x)]}}
\RightLabel{\TRule{\True}{\forall}}
\UnaryInfC{\sFmla{\True}{A(t)}}
\DisplayProof
&
\AxiomC{\sFmla{\False}{\lforall[x][A(x)]}}
\RightLabel{\TRule{\False}{\lforall}}
\UnaryInfC{\sFmla{\False}{A(a)}}
\DisplayProof
\\[6ex]
\AxiomC{\sFmla{\True}{\lexists[x][A(x)]}}
\RightLabel{\TRule{\True}{\lexists}}
\UnaryInfC{\sFmla{\True}{A(a)}}
\DisplayProof
&
\AxiomC{\sFmla{\False}{\lexists[x][A(x)]}}
\RightLabel{\TRule{\False}{\lexists}}
\UnaryInfC{\sFmla{\False}{A(t)}}
\DisplayProof
\end{tabular}
\end{center}
\caption{\Log{Tc} चे नियम. $\TRule{\False}{\lforall}$ आणि
$\TRule{\True}{\lexists}$ मध्ये स्थिरांक~$a$ वरच्या
शाखेत आलेला नसावा; म्हणजे तो शाखेसाठी नवा असावा.}
\label{pt:ps:tab:tab:Tc}
\end{table}




नियम म्हणजे~\Log{G3c} चेच नियम उलटे
दाखवले आहेत. चिन्हे~$\True$ आणि~$\False$
अनुक्रमे क्रमवर्तीच्या डाव्या किंवा उजव्या बाजूला
सूत्र मुख्य आहे हे दर्शवतात.

टॅब्लोची शाखा~$\sFmla{\True}{A}$ आणि
$\sFmla{\False}{A}$ ही दोन्ही सूत्रे धरत असेल,
किंवा तिच्यावर~$\sFmla{\True}{\lfalse}$ असेल,
तर ती \emph{बंद} आहे. प्रत्येक शाखा बंद असेल,
तर संपूर्ण टॅब्लो बंद म्हणतो. उदाहरणार्थ,
$A \lor (B \land C) \Sequent (A \lor B) \land
(A \lor C)$ साठी हा बंद टॅब्लो आहे; म्हणजे
$\sFmla{\True}{A \lor (B \land C)}$ आणि
$\sFmla{\False}{(A \lor B) \land (A \lor C)}$
यांपासून सुरुवात होते:
\begin{oltableau}
  [\sFmla{\True}{\formula{A} \lor (\formula{B} \land \formula{C})},
    just=\TAss, checked
    [\sFmla{\False}{(\formula{A} \lor \formula{B}) \land
        (\formula{A} \lor \formula{C})}, just=\TAss, checked
      [\sFmla{\True}{\formula{A}}, just={\TRule{\True}{\lor}[1]}
        [\sFmla{\False}{\formula{A} \lor \formula{B}},
          just={\TRule{\False}{\land}[2]}, checked
          [\sFmla{\False}{\formula{A}},
            just={\TRule{\False}{\lor}[4]}
            [\sFmla{\False}{\formula{B}},
              just={\TRule{\False}{\lor}[4]}, close]
          ]
        ]
        [\sFmla{\False}{\formula{A} \lor \formula{C}},
          just={\TRule{\False}{\land}[2]}, checked
          [\sFmla{\False}{\formula{A}},
            just={\TRule{\False}{\lor}[4]}
            [\sFmla{\False}{\formula{C}},
              just={\TRule{\False}{\lor}[4]}, close]
          ]
        ]
      ]
      [\sFmla{\True}{\formula{B} \land \formula{C}},
        just={\TRule{\True}{\lor}[1]}, checked
        [\sFmla{\False}{\formula{A} \lor \formula{B}},
          just={\TRule{\False}{\land}[2]}, checked
          [\sFmla{\True}{\formula{B}},
            just={\TRule{\True}{\land}[3]}, move by=2
            [\sFmla{\True}{\formula{C}},
              just={\TRule{\True}{\land}[3]}
              [\sFmla{\False}{\formula{A}},
                just={\TRule{\False}{\lor}[4]}
                [\sFmla{\False}{\formula{B}},
                  just={\TRule{\False}{\lor}[4]},close
                ]
              ]
            ]
          ]
        ]
        [\sFmla{\False}{\formula{A} \lor \formula{C}},
          just={\TRule{\False}{\land}[2]}, checked
          [\sFmla{\True}{\formula{B}},
            just={\TRule{\True}{\land}[3]}, move by=2
              [\sFmla{\True}{\formula{C}},
                just={\TRule{\True}{\land}[3]}
                [\sFmla{\False}{\formula{A}},
                  just={\TRule{\False}{\lor}[4]}
                  [\sFmla{\False}{\formula{C}},
                  just={\TRule{\False}{\lor}[4]},close
                ]
              ]
            ]
          ]
        ]
      ]
    ]
  ]
\end{oltableau}

खूणा दर्शवतात की संबंधित नियम त्या
सूत्राच्या खालच्या सर्व शाखांवर लागू
केला आहे. ~$\otimes$ हे चिन्ह शाखा बंद असल्याचे
दर्शवते. नियमानंतरच्या ओळ-क्रमांकांवरून नियम
कोणत्या चिन्हांकित सूत्रे ला लागू केला
हे कळते (म्हणजे दर्शवलेल्या ओळीवरील त्याच
शाखेतील सूत्राला).

टॅब्लोमध्ये सिद्धता-शोध टप्प्याटप्प्याने करतो.
येथील बंद-पद शोधासाठी सुरुवातीची सूत्रे वाक्ये असावीत.
$0$ व्या टप्प्यावर ती वर लिहा. प्रत्येक टप्प्यावर
बंद शाखा पुढे वाढवू नका. प्रत्येक उघड्या शाखेत अद्याप
खूण न केलेले सर्वांत वरचे अणुरूप नसलेले सूत्र निवडा,
पण पुढील परिच्छेदातील पुन्हा वापरायचे संख्यापक त्यातून
वगळा. त्याला शाखा-विस्तार नियम लावा. संबंधित सर्व
उघड्या शाखांवर नियम लागू झाल्यावर त्याला खूण करा.
अणुरूप सूत्रांना विस्तार-नियम लागत नाही.
(वरचे विधानीय उदाहरण या पद्धतीने बनवले आहे.)

$\sFmla{\True}{\lforall}$ आणि
$\sFmla{\False}{\lexists}$ ही येथे संख्यापक-प्रकारांची
संक्षिप्त रूपे आहेत; संख्यापक एकटाच पूर्ण सूत्र असल्याचा
अर्थ नाही. पूर्ण रूपे~$\sFmla{\True}{\lforall[x][B(x)]}$
आणि~$\sFmla{\False}{\lexists[x][B(x)]}$ अशी आहेत.
या सूत्रांवर कायमची खूण करत नाही: सर्व योग्य पदांसाठी
अनुरूप नियम अखेरीस लावले गेले पाहिजेत.
प्रत्येक उघड्या शाखेबरोबर स्थिरांकांचा संच~$C$ ठेवा.
सुरुवातीला त्यात सुरुवातीच्या सूत्रांतील सर्व स्थिरांक
घ्या; ते नसतील, तर~$\{\Obj c_0\}$ घ्या.
$\TRule{\False}{\lforall}$ किंवा~$\TRule{\True}{\lexists}$
साठी आणलेला नवा स्थिरांक~$C$ मध्ये जोडा, तो सूत्रात
प्रत्यक्ष उरला नाही तरी; तो आधीच्या~$C$ बाहेरचा असावा.
शाखा फुटल्यावर दोन्ही शाखांना त्या संचाची प्रत द्या.
सामान्य सूत्र हाताळल्यानंतर, त्या टप्प्यावर शाखेतील
प्रत्येक पुन्हा वापरायचा संख्यापक एकदा हाताळा.
निश्चित पद-क्रमातून~$C$ मधील स्थिरांक आणि सुरुवातीची
फलनचिन्हे वापरून बनणारे पहिले पद~$t$ निवडा ज्यासाठी
अनुरूप~$\sFmla{\True}{B(t)}$ किंवा
$\sFmla{\False}{B(t)}$ शाखेत नाही, आणि ते उदाहरण जोडा.
असे पद नसेल, तर त्या टप्प्यावर या सूत्रासाठी काही
जोडू नका; पुढे~$C$ वाढल्यास पुन्हा तपासा.

या शोधातून बंद टॅब्लो मिळाला नाही, तर सांत उघडी
शाखा पूर्णपणे संतृप्त झाली असेल किंवा शोध अनंत,
सांत-शाखित वृक्ष तयार करत राहील. येथे सांत शाखा
संतृप्त असण्याचा अर्थ असा: पुन्हा वापरायचे संख्यापक
सोडून सर्व अणुरूप नसलेली सूत्रे हाताळली आहेत,
आणि प्रत्येक पुन्हा वापरायच्या संख्यापकासाठी~$C$ व
सुरुवातीच्या फलनचिन्हांपासून बनणाऱ्या प्रत्येक पदाचे
अनुरूप चिन्हांकित उदाहरण आधीच शाखेत आहे.
अनंत वृक्षाला अनंत शाखा असते. निश्चित न्याय्य पद-क्रमामुळे
तिच्यावरही प्रत्येक आवश्यक संख्यापक-उदाहरण अखेरीस येते.
\ref{pt:ps:cpl:sec} प्रमाणे~$\Struct{M}$ बांधता येते:
$\Theta = \Setabs{A}{\sFmla{\True}{A} \text{ शाखेवर}}$
आणि~$\Xi = \Setabs{A}{\sFmla{\False}{A} \text{ शाखेवर}}$
घ्या; स्थिरांकांसाठी शाखेबरोबर नोंदलेल्या संचांचे संघटन
घ्या. तो संच अरिक्त आहे. मग शाखेवर
$\sFmla{\True}{A}$ असेल तर~$\Sat{M}{A}$,
आणि~$\sFmla{\False}{A}$ असेल तर~$\Sat/{M}{A}$.
संतृप्त शाखेसाठी संख्यापक-उदाहरणे सर्व पदांसाठी उपलब्ध
असल्याने पद-प्रतिमानातील पूर्वीचे विगमन लागू होते.

या पद्धतीने तयार झालेले टॅब्लो प्रत्येक
शिखराला क्रमवर्ती देऊन~\Log{G3c} मधील सिद्धतांमध्ये
सहज रूपांतरित करता येतात. सुरुवातीची
सूत्रे ही~$\Gamma \Sequent \Delta$
क्रमवर्तीशी संबंधित असतील, तर त्या प्रत्येक
सूत्राला~$\Gamma \Sequent \Delta$ ही खूण द्या.
~$\sFmla{\True}{A}$ या चिन्हांकित सूत्र
वर शाखा-विस्तार नियम लावून शाखा वाढवली, तर
पानाला~$A, \Pi \Sequent \Lambda$ ही खूण द्या;
~$\sFmla{\False}{A}$ या चिन्हांकित सूत्र
वर नियम लावला,
तर पानाला~$\Pi \Sequent \Lambda, A$ ही खूण द्या.
टॅब्लोत~\TRule{\True}{\star} नियम वापरला,
तर क्रमवर्तीवर~\LeftR{\star} नियम लावा
(मुख्य सूत्र~$A$);~\TRule{\False}{\star}
वापरला असेल, तर~\RightR{\star} लावा.
नियमाची आधारविधाने ही टॅब्लोत नव्याने
जोडलेल्या अनुरूप शिखरांना दिलेल्या क्रमवर्ती आहेत.
~$\sFmla{\True}{A}$ आणि~$\sFmla{\False}{A}$
यांनी टॅब्लो बंद झाल्यावर अंतिम पानशिखराची
खूण~$A, \Pi \Sequent \Lambda, A$
या रूपातील क्रमवर्ती असेल. $\sFmla{\True}{\lfalse}$
ने शाखा बंद झाली असेल, तर पानाची खूण
$\lfalse,\Pi\Sequent\Lambda$ या स्वयंसिद्धकाचे रूप असेल.
टॅब्लोतील काही सलग
शिखरांना तीच क्रमवर्ती दिली जाईल; तिच्या
पुनरावृत्ती काढा. उदाहरणार्थ, वरच्या टॅब्लोचे
पुढील सिद्धता मध्ये रूपांतर होते:
\[\small
\Axiom$A \fCenter A, B$
\RightLabel{\RightR{\lor}}
\UnaryInf$A \fCenter A \lor B$
\Axiom$A \fCenter A, C$
\RightLabel{\RightR{\lor}}
\UnaryInf$A \fCenter A \lor C$
\RightLabel{\RightR{\land}}
\BinaryInf$A \fCenter (A \lor B) \land (A \lor C)$
\Axiom$B, C \fCenter A, B$
\RightLabel{\RightR{\lor}}
\UnaryInf$B, C \fCenter A \lor B$
\RightLabel{\LeftR{\land}}
\UnaryInf$B \land C \fCenter A \lor B$
\Axiom$B, C \fCenter A, C$
\RightLabel{\RightR{\lor}}
\UnaryInf$B, C \fCenter A \lor C$
\RightLabel{\LeftR{\land}}
\UnaryInf$B \land C \fCenter A \lor C$
\RightLabel{\RightR{\land}}
\BinaryInf$B \land C \fCenter (A \lor B) \land
(A \lor C)$
\RightLabel{\LeftR{\lor}}
\BinaryInf$A \lor (B \land C) \fCenter (A \lor B) \land
(A \lor C)$
\DisplayProof
\]



